%% Technical Chapters

\chapter{System Overview}

The EvoRAG system is designed as a modular, privacy-preserving pipeline that enhances
traditional Retrieval-Augmented Generation by incorporating multi-agent reasoning,
evolutionary learning, and user-centric adaptation.

\section{Functional and Non-Functional Requirements}

\subsection{Functional Requirements}

\begin{itemize}

    \item \textbf{Real-Time Knowledge Ingestion:}
    The system shall fetch live article content at query time from NewsAPI, GNews,
    and a curated set of RSS feeds (BBC News, TechCrunch, Reuters) using the
    \texttt{requests} and \texttt{feedparser} libraries, triggered exclusively by
    the user's submitted query keyword — no pre-stored or pre-fetched corpus shall
    be used.

    \item \textbf{HTML Cleaning and Chunk Segmentation:}
    The system shall strip all non-content elements (CSS, JavaScript, navigation,
    advertisements) from fetched web pages using BeautifulSoup4, deduplicate articles
    by SHA-256 URL hashing, discard articles below a minimum character threshold,
    and segment retained content into overlapping chunks of 300--400 words with a
    50-word overlap to preserve cross-boundary context.

    \item \textbf{Vector Embedding and FAISS Indexing:}
    The system shall encode every text chunk into a dense vector using the
    \texttt{all-MiniLM-L6-v2} sentence-transformer model and persist all embeddings
    in a session-scoped local FAISS index alongside metadata (source URL, title,
    publication timestamp) to support semantic nearest-neighbour retrieval.

    \item \textbf{Contextual Chunk Retrieval:}
    The system shall, for each incoming user query, embed the query using the same
    \texttt{all-MiniLM-L6-v2} model and retrieve the top-3 semantically closest
    chunks from the FAISS index via approximate nearest-neighbour search, forming
    the grounding context for all subsequent persona generation calls.

    \item \textbf{Eight-Persona Parallel Response Generation:}
    The system shall invoke a locally hosted Ollama LLM (Llama~3.1 or Mistral~7B)
    once per query cycle with a structured prompt that instructs the model to
    generate eight independent analytical responses in a single call, each
    representing a distinct reasoning style: Analytical-Critical, Empirical-Evidential,
    Adversarial-Sceptical, Creative-Lateral, Domain-Expert, Temporal-Contextual,
    Risk-Evaluative, and Synthesis-Integrative.

    \item \textbf{Automated Peer-Review Scoring:}
    The system shall pass all eight persona drafts to the same local LLM in a
    second structured prompt, instructing it to act as an impartial judge and assign
    each response a score out of 30 across three dimensions — Factual Grounding (0--10),
    Analytical Depth (0--10), and Clarity (0--10) — returning results as a
    parseable JSON object.

    \item \textbf{Consensus-Based Answer Selection:}
    The system shall parse the judge's JSON score output, aggregate total scores per
    persona, select the persona with the highest aggregate score as the winner, and
    surface that response as the final answer together with the complete numerical
    score breakdown for all eight personas.

    \item \textbf{Source Attribution in Final Output:}
    The system shall append to every final answer the source URL, article title, and
    publication timestamp of each of the top-3 retrieved chunks that were used as
    the grounding context for that query cycle.

    \item \textbf{Evolutionary Win-Rate Tracking:}
    After every query cycle, the system shall update a local SQLite database with
    the winning persona's identifier, its total score, and any explicit user feedback
    signal (thumbs-up or thumbs-down), accumulating a longitudinal performance record
    per persona across all sessions.

    \item \textbf{Adaptive Persona Replacement:}
    The system shall compute a rolling win-rate for each persona from the SQLite
    history and, when a persona's win-rate falls below a configurable threshold,
    replace its prompt configuration with a newly generated variant without modifying
    the underlying Ollama model weights or restarting the backend service.

    \item \textbf{User Feedback and Voting Weight Adjustment:}
    The system shall record per-query user feedback via the React frontend and use
    accumulated feedback to adjust the relative voting weight assigned to each persona
    in the consensus scoring phase, such that personas consistently preferred by a
    specific user receive a weighted scoring advantage in future queries.

\end{itemize}

\subsection{Non-Functional Requirements}

\begin{itemize}

    \item \textbf{Local-Only Execution:}
    Every component of EvoRAG — including the Ollama LLM runtime, FAISS vector index,
    SQLite persona database, FastAPI backend, and React frontend — shall run entirely
    on the user's local machine. The system shall make no outbound network calls during
    the generation, scoring, or evolutionary update phases; network access is permitted
    only during the initial knowledge acquisition phase (NewsAPI, GNews, RSS).

    \item \textbf{Retrieval-Grounded Generation:}
    The LLM prompt issued for persona generation shall always include the top-3
    retrieved FAISS chunks as mandatory context. The system shall not issue any
    generation call without valid retrieval context; if the FAISS index returns zero
    results, the pipeline shall terminate early and return a structured
    ``no relevant context'' message rather than invoking the LLM on an empty context.

    \item \textbf{End-to-End Response Latency:}
    On the minimum hardware configuration (CPU-only inference with Llama~3.1 8B or
    Mistral~7B via Ollama), the complete pipeline — fetch, embed, retrieve, generate
    eight personas, score, and select — shall complete within 90 seconds per query.
    On the recommended configuration (NVIDIA RTX 3060 or equivalent with CUDA-enabled
    Ollama), this target reduces to 25 seconds per query.

    \item \textbf{FAISS Index Capacity and Retrieval Consistency:}
    The FAISS flat-L2 index shall handle a per-session knowledge base of up to
    10,000 text chunks (approximately 50--100 fully processed articles) with no
    measurable increase in nearest-neighbour search latency, as FAISS flat indices
    perform exact search in O(n) time over the embedding dimension, providing
    deterministic retrieval results for identical query embeddings.

    \item \textbf{Graceful Handling of News API Failures:}
    If NewsAPI or GNews returns a non-200 HTTP status or a response with zero
    articles, the pipeline shall automatically fall back to the configured RSS feed
    set and proceed with whatever content is available. If all three sources yield
    no usable content, the system shall return a structured failure response; it
    shall not raise an unhandled exception or crash the FastAPI server process.

    \item \textbf{Persona Score Transparency:}
    The React frontend shall display the full scoring table — all eight persona names,
    their individual Factual Grounding, Analytical Depth, and Clarity scores, and
    aggregate totals — alongside the winning response for every completed query.
    This display shall not be optional or hidden behind additional user interaction.

    \item \textbf{SQLite Schema Stability Across Sessions:}
    The evolutionary tracking database (SQLite) shall persist across server restarts.
    The schema (persona identifier, win count, total queries, last updated timestamp,
    user feedback score) shall not be modified between sessions; accumulated win-rate
    history shall carry forward and inform persona replacement decisions continuously.

\end{itemize}

\section{Software Requirements}

The EvoRAG system relies entirely on software-based components and locally hosted
models, enabling full execution without cloud dependencies or external API calls.
Table~\ref{tab:software} summarises the complete software environment.

\vspace{0.5em}

\begin{center}
\small
\renewcommand{\arraystretch}{1.5}
\begin{longtable}{|p{3.8cm}|p{3.5cm}|p{6.2cm}|}
\caption{Software Requirements}
\label{tab:software} \\
\hline
\textbf{Component} & \textbf{Specification} & \textbf{Purpose} \\
\hline
\endfirsthead
\hline
\textbf{Component} & \textbf{Specification} & \textbf{Purpose} \\
\hline
\endhead
\hline
\endfoot
Programming Language & Python 3.11+ &
Extensive ML ecosystem, cross-platform compatibility, and asynchronous
programming support for the multi-persona generation pipeline. \\
\hline
Embedding Library & sentence-transformers $\geq$ 3.1.0 &
Converts text chunks into high-dimensional vector representations for
semantic retrieval using the all-MiniLM-L6-v2 model. \\
\hline
Vector Search Library & faiss-cpu $\geq$ 1.9.0 &
Provides efficient vector indexing and approximate nearest-neighbour
lookup for the retrieval pipeline. \\
\hline
Numerical Computing & numpy $\geq$ 1.26 &
Enables efficient array operations and mathematical computations
throughout the pipeline. \\
\hline
Web Framework & fastapi $\geq$ 0.115.0 &
Exposes the EvoRAG pipeline as a RESTful backend API with automatic
OpenAPI documentation. \\
\hline
ASGI Server & uvicorn[standard] $\geq$ 0.30.0 &
Serves the FastAPI application with support for asynchronous request
handling and HTTP/2. \\
\hline
HTTP Client & httpx $\geq$ 0.27.0 &
Handles asynchronous inter-service communication within the pipeline. \\
\hline
Feed Parser & feedparser $\geq$ 6.0.11 &
Parses RSS and Atom feeds during the knowledge acquisition phase. \\
\hline
HTML Parser & beautifulsoup4 $\geq$ 4.12.3 &
Cleans and extracts textual content from fetched web articles. \\
\hline
HTTP Library & requests $\geq$ 2.32.0 &
Supports synchronous web requests during the data fetching phase. \\
\hline
Environment Manager & python-dotenv $\geq$ 1.0.1 &
Manages environment variables and configuration securely at runtime. \\
\hline
Local LLM Runtime & Ollama $\geq$ 0.3.0 with Llama~3.1~(8B) / Mistral~7B &
Hosts and serves open-source large language models entirely on local
hardware, eliminating any third-party API dependency. \\
\hline
Embedding Model & nomic-embed-text (via Ollama) &
Provides a locally hosted text embedding model as an alternative to
sentence-transformers for fully offline operation. \\
\hline
Vector Store & FAISS (local, in-process) &
Persists and queries embedded knowledge chunks with no external
database server required. \\
\hline
Database & SQLite 3 (built-in) &
Stores persona performance metrics, win rates, and evolutionary
history in a lightweight, file-based local database. \\
\hline
Frontend & React.js $\geq$ 18.0 &
Provides real-time query input, persona draft display, voting breakdown
visualisation, and user feedback controls. \\
\hline
Operating System & Windows 10/11 or Ubuntu 22.04+ &
Required to support the Python 3.11+ runtime and the Ollama local
model serving platform. \\
\hline
Development Environment & Visual Studio Code $\geq$ 1.89 &
Offers integrated debugging, code completion, and version control
capabilities. \\
\hline
\end{longtable}
\end{center}

\clearpage

\section{Hardware Requirements}

The EvoRAG system runs entirely on a single local development workstation.
Table~\ref{tab:hardware} summarises the minimum and recommended hardware configuration.

\vspace{0.5em}

\begin{table}[H]
\centering
\caption{Hardware Requirements}
\label{tab:hardware}
\renewcommand{\arraystretch}{1.5}
\footnotesize
\begin{tabularx}{\linewidth}{|p{2.2cm}|X|X|X|}
\hline
\textbf{Component} & \textbf{Minimum} & \textbf{Recommended} & \textbf{Justification} \\
\hline
Processor &
Intel Core i5 (8th Gen) / AMD Ryzen 5 3600 &
Intel Core i7 (12th Gen) / AMD Ryzen 7 5800X &
Required for local LLM inference and concurrent multi-persona generation. \\
\hline
RAM &
16 GB DDR4 &
32 GB DDR4/DDR5 &
Supports simultaneous execution of Ollama LLM, FAISS search, FastAPI backend, and React frontend. \\
\hline
GPU &
NVIDIA GTX 1660 (6 GB VRAM) or CPU-only &
NVIDIA RTX 3060 (12 GB VRAM) or RTX 4070 &
Accelerates LLM inference via CUDA; system runs on CPU by default but GPU reduces latency significantly. \\
\hline
Storage &
20 GB SSD &
50 GB NVMe SSD &
Accommodates LLM weights (4--8 GB), FAISS index, SQLite database, and fetched knowledge articles. \\
\hline
Operating System &
Windows 10 (64-bit) or Ubuntu 20.04 LTS &
Windows 11 or Ubuntu 22.04 LTS &
Required for Python 3.11+ runtime and Ollama local model serving platform. \\
\hline
Network &
\multicolumn{2}{X|}{Broadband internet (knowledge acquisition phase only)} &
Required only during initial fetch via NewsAPI/GNews/RSS; all inference and retrieval run fully offline. \\
\hline
\end{tabularx}
\end{table}

\section{System Architecture}

The EvoRAG architecture is structured as a continuous, evolutionary loop consisting of ten primary stages. This architecture ensures that the system not only provides grounded answers but also improves its analytical capabilities over time based on performance metrics and user feedback.

The workflow follows these stages:

\begin{enumerate}
    \item \textbf{User Query:} The process begins when a user submits a question to the system.
    
    \item \textbf{Data Retrieval:} Relevant documents and raw data are fetched from various APIs and external sources.
    
    \item \textbf{Processing \& Embedding:} The retrieved data is cleaned, segmented into chunks, and converted into high-dimensional vector embeddings.
    
    \item \textbf{Vector Search:} The system searches the vector database to retrieve the most contextually relevant chunks for the specific query.
    
    \item \textbf{Multi-Persona Generation:} Responses are generated in parallel by multiple distinct AI personas (e.g., Analyst, Critic, Creative).
    
    \item \textbf{Consensus Voting:} The personas evaluate each other’s responses, scoring them based on quality, factual accuracy, and depth.
    
    \item \textbf{Final Answer:} The highest-scoring response is selected and delivered as the final output to the user.
    
    \item \textbf{User Feedback:} The user provides qualitative feedback (likes or dislikes), which is recorded by the system.
    
    \item \textbf{Performance Tracking:} The system tracks persona accuracy and scores, storing these performance metrics in a local database.
    
    \item \textbf{Evolutionary Update:} Based on performance trends, weak personas are replaced or refined, allowing the system to improve autonomously over time.
\end{enumerate}

\begin{figure}[H]
    \centering
    \includegraphics[width=0.70\linewidth]{Sys.png}
    \caption{EvoRAG System Flowchart}
    \label{fig:evorag_sys_flowchart}
\end{figure}

\begin{figure}[H]
    \centering
    \includegraphics[width=0.95\linewidth]{sA2.png}
    \caption{EvoRAG System Architecture and Evolutionary Loop}
    \label{fig:evorag_architecture}
\end{figure}


\subsection{High Level Design}

The high level design illustrates the major modules of the EvoRAG system and their interactions.

\begin{figure}[H]
    \centering
    \includegraphics[width=0.85\linewidth,height=0.85\textheight,keepaspectratio]{hld.png}
    \caption{EvoRAG High Level Design}
    \label{fig:hld}
\end{figure}

\clearpage

\subsection{Low Level Design}

The low level design provides a detailed view of the internal logic within each
module, specifying the data flows, component interactions, API contracts, and
the sequence of operations that occur during a complete query-to-response cycle,
including the evolutionary feedback loop.

\begin{figure}[H]
    \centering
    \includegraphics[width=\linewidth,height=1.0\textheight,keepaspectratio]{lld3.jpg}
    \caption{EvoRAG Low Level Design}
    \label{fig:lld}
\end{figure}

\chapter{Methodology}

\begin{figure}[H]
    \centering
    \includegraphics[width=0.75\linewidth]{method.png}
    \caption{Proposed Methodology}
    \label{fig:method_flowchart}
\end{figure}

Figure~\ref{fig:method_flowchart} illustrates the complete EvoRAG methodology flow. The methodology of EvoRAG is divided into four core engines that work in tandem to produce high-quality, personalized results. It begins with dynamic data acquisition, progresses through data cleaning and chunking, indexes the resulting semantic vectors, and finishes with parallel multi-persona reasoning evaluated through consensus.

\section{Implementation}
\label{sec:implementation}

This section details the physical realization of the EvoRAG pipeline, translating the four core architectural engines into an end-to-end, locally executable system. The implementation encompasses live multi-source knowledge acquisition, text sanitization and deduplication, sliding-window chunk segmentation, dense vector embedding generation in FAISS, semantic retrieval with dynamic synchronization, parallel multi-persona generation across asynchronous semaphores, background scheduler orchestration, and browser-native voice and chat interfaces.

\subsection{Implementation Overview}
\label{subsec:impl_overview}
The physical implementation is structured across four interconnected architectural tiers executing within an asynchronous runtime environment:
\begin{enumerate}[label=(\roman*)]
    \item \textbf{Data Ingestion and Indexing Tier:} Orchestrates scheduled and on-demand news harvesting across NewsAPI, GNews, and curated RSS feeds, DOM extraction, HTML boilerplate stripping, SHA-256 URL and text deduplication, sliding-window chunking, and incremental vectorization using dense Sentence-Transformers persisted in a flat FAISS index.
    \item \textbf{Semantic Retrieval and Synchronization Tier:} Performs nearest-neighbour vector search over indexed chunk embeddings, evaluates context sufficiency and temporal freshness, and dynamically reloads updated index vectors in memory via timestamp tracking without server restarts.
    \item \textbf{Multi-Persona Reasoning and Consensus Tier:} Orchestrates eight specialized analytical personas executing concurrently across an asynchronous semaphore pool, facilitates multi-criteria peer evaluation, and synthesizes the winning arguments into a unified consensus output.
    \item \textbf{Service, Persistence, and Presentation Tier:} Exposes asynchronous REST endpoints via FastAPI, persists multi-turn conversation sessions and user preferences in SQLite via SQLModel, manages background task scheduling via APScheduler, and renders an interactive web interface supporting browser-native speech-to-text (STT) and text-to-speech (TTS) modalities.
\end{enumerate}

\subsection{Dataset and Knowledge Acquisition Pipeline}
\label{subsec:data_acquisition}

\subsubsection{Live Multi-Source Corpus Harvesting}
Unlike static open-domain question answering systems that rely on pre-compiled, frozen corpora such as Wikipedia dumps, EvoRAG constructs its runtime knowledge base dynamically from live, multi-source digital news feeds. This design prevents temporal obsolescence and enables the system to analyze breaking events and technical developments. Data acquisition is orchestrated by \texttt{fetcher.py} using three heterogeneous channels:
\begin{itemize}
    \item \textbf{NewsAPI REST Endpoint:} Queries global journalistic publishers by keyword, specifying English-language constraints, recency sorting (\texttt{sortBy=publishedAt}), and page size limits.
    \item \textbf{Google News (GNews) API:} Accesses aggregated Google News indexing to capture top-ranking mainstream reporting across regional editions.
    \item \textbf{Curated Technical RSS Feeds:} Continuously monitors high-credibility syndicate feeds using the \texttt{feedparser} library. Monitored sources include BBC World News, Reuters Technology, MIT Technology Review (AI), The Verge, VentureBeat AI, Google AI Blog, and OpenAI News.
\end{itemize}

Table~\ref{tab:news_sources} outlines the ingestion sources and operational characteristics configured in \texttt{config.py}.

\begin{table}[H]
\centering
\caption{Knowledge Ingestion Source Configuration}
\label{tab:news_sources}
\small
\renewcommand{\arraystretch}{1.3}
\begin{tabularx}{\linewidth}{|l|p{3.2cm}|p{2.8cm}|X|}
\hline
\textbf{Source Channel} & \textbf{Protocol / Format} & \textbf{Query Mechanism} & \textbf{Coverage Scope} \\
\hline
NewsAPI & HTTPS REST (JSON) & Keyword search endpoint & Global mainstream news publishers \\
\hline
GNews API & HTTPS REST (JSON) & Topic search endpoint & Top-ranked international news articles \\
\hline
BBC World News & RSS 2.0 (XML) & Syndicated feed polling & International geopolitical and world reporting \\
\hline
Reuters Tech & RSS 2.0 (XML) & Syndicated feed polling & Global technology and market developments \\
\hline
MIT Tech Review & RSS 2.0 / Atom (XML) & Topic-filtered feed & Artificial intelligence and applied science research \\
\hline
The Verge / VentureBeat & RSS 2.0 (XML) & Category feed polling & Industry computing, semiconductor and enterprise news \\
\hline
Google AI / OpenAI Blog & Atom / RSS (XML) & Official feed polling & Frontier model research and alignment announcements \\
\hline
\end{tabularx}
\end{table}

\subsubsection{Data Cleaning and Noise Removal}
Raw web documents extracted from digital news feeds contain substantial non-informative structural overhead, including navigation menus, inline scripts, styling directives, tracking pixels, and advertising banners. To prevent context window pollution and maintain high signal-to-noise ratios, incoming HTML responses are processed through a deterministic multi-stage cleaning pipeline:
\begin{enumerate}
    \item \textbf{DOM Parsing and Tree Sanitization:} The document body is parsed using \texttt{BeautifulSoup4} with the native \texttt{lxml} or \texttt{html.parser} engine. Elements tagged as \texttt{<script>}, \texttt{<style>}, \texttt{<nav>}, \texttt{<header>}, \texttt{<footer>}, \texttt{<aside>}, \texttt{<form>}, \texttt{<noscript>}, or \texttt{<iframe>} are recursively decomposed.
    \item \textbf{Semantic Body Extraction:} The parser prioritizes semantic container tags (\texttt{<article>}, \texttt{<main>}) before falling back to \texttt{<body>}, extracting continuous narrative text while filtering navigation lists.
    \item \textbf{Whitespace and Encoding Normalization:} Multiple consecutive whitespace characters, line breaks, and HTML entity encodings are normalized into standard UTF-8 whitespace-separated tokens.
    \item \textbf{Length Filtering:} Articles yielding fewer than 100 words following sanitization are discarded, eliminating stub pages, redirection notices, and paywalled previews.
\end{enumerate}

\subsubsection{Deduplication Strategy}
To prevent duplicate reporting from flooding the retrieval index when multiple outlets publish syndicated coverage of the same event, deduplication is enforced at both the document and chunk levels:
\begin{equation}
\label{eq:url_hash}
h_{\text{url}} = \text{MD5}(\text{normalize\_url}(U))
\end{equation}
where $\text{normalize\_url}(U)$ strips trailing slashes, URL tracking parameters (\texttt{utm\_*}, \texttt{ref}), and protocol variations. If $h_{\text{url}}$ matches an existing article identifier in the current index or session database, the document is discarded prior to chunking. Furthermore, chunk-level deduplication generates a SHA-256 cryptographic digest of the normalized chunk body text to detect duplicate paragraphs across distinct syndication URLs:
\begin{equation}
\label{eq:text_hash}
h_{\text{text}} = \text{SHA-256}(\text{strip}(\text{text}))
\end{equation}

\subsubsection{Sliding-Window Chunking Strategy}
Large Language Models operate under strict context window and token limitations, making it impractical to feed entire articles as input. To address this, a character-bounded sliding-window chunking strategy is applied.

Let an article's word sequence be $W = (w_1, w_2, \dots, w_M)$. The sequence is partitioned into overlapping subsets:
\begin{equation}
\label{eq:chunk_window}
C_k = (w_{s_k}, w_{s_k + 1}, \dots, w_{e_k})
\end{equation}
where the stride step is governed by chunk target length $L_{\text{target}} = 350$ words and overlap length $L_{\text{overlap}} = 50$ words:
\begin{equation}
s_k = k \cdot (L_{\text{target}} - L_{\text{overlap}}) + 1, \quad e_k = \min(s_k + L_{\text{target}} - 1, M)
\end{equation}
Chunks shorter than $L_{\text{min}} = 300$ words are discarded unless they represent the terminal segment of an article exceeding the minimum threshold. Chunks exceeding $L_{\text{max}} = 400$ words are split recursively. Each chunk is packaged with a metadata dictionary preserving the source publication name, canonical article URL, publication timestamp, chunk ordinal index, and document title:
\begin{equation}
\mathcal{M}_k = \{ \text{chunk\_id}, \text{source}, \text{url}, \text{url\_hash}, \text{title}, \text{published\_at}, \text{chunk\_index}, \text{total\_chunks}, \text{word\_count}, \text{text} \}
\end{equation}

\subsection{Models, Embedding Representations, and Vector Database}
\label{subsec:models_faiss}

\subsubsection{Dense Vector Embedding Generation}
Semantic vectorization is implemented in \texttt{embedder.py} utilizing the \texttt{sentence-transformers} library. The system employs \texttt{all-MiniLM-L6-v2}, a distilled 6-layer bidirectional transformer pre-trained on over one billion sentence pairs. 

Given a sanitized text passage $C_k$, the model computes token-level contextual representations and applies mean pooling over the output layer to yield a dense vector $\mathbf{v}_k \in \mathbb{R}^{D}$ where $D = 384$:
\begin{equation}
\label{eq:mean_pooling}
\mathbf{v}_k = \frac{\sum_{i=1}^{T_k} m_i \mathbf{h}_i}{\sum_{i=1}^{T_k} m_i}
\end{equation}
where $\mathbf{h}_i$ is the $i$-th hidden token vector, $m_i \in \{0, 1\}$ is the attention mask, and $T_k$ is the sequence length. Vectorization is performed in mini-batches of size $B = 64$ with PyTorch running in CPU-optimized FP32 precision.

\subsubsection{FAISS Vector Store Architecture}
Nearest-neighbour similarity search is supported via Facebook AI Similarity Search (FAISS). EvoRAG adopts an exact Euclidean flat index structure (\texttt{IndexFlatL2}). For query embedding $\mathbf{q} \in \mathbb{R}^{D}$ and indexed passage vectors $\mathbf{v}_k \in \mathbb{R}^{D}$, the retrieval distance is:
\begin{equation}
\label{eq:l2_distance}
d(\mathbf{q}, \mathbf{v}_k) = \|\mathbf{q} - \mathbf{v}_k\|_2 = \sqrt{\sum_{j=1}^{D} (q_j - v_{k,j})^2}
\end{equation}
Because the MiniLM-L6-v2 embeddings operate within a compact 384-dimensional manifold, \texttt{IndexFlatL2} provides deterministic, exact retrieval with sub-millisecond query latency on local workstations for corpus sizes up to $50,000$ chunks, eliminating the quantization noise and hyperparameter tuning required by approximate inverted index schemes (IVF-PQ).

Alongside the binary index file (\texttt{index/evorag.faiss}), a companion JSON sidecar (\texttt{index/evorag\_meta.json}) stores aligned metadata records indexed by their contiguous row identifiers.

\subsubsection{Incremental Vector Updates and Atomic Persistence}
In dynamic news environments, rebuilding the entire vector index from scratch upon each news ingestion cycle is computationally wasteful. In \texttt{embedder.py}, the \texttt{add\_chunks\_to\_index()} function implements incremental vector addition. Newly discovered chunks are identified by comparing their cryptographic URL-index and text signatures against the existing metadata set. Only genuine novel chunks are passed to the embedding model.

The resulting embeddings are directly appended to the live index:
\begin{equation}
\text{index.add}(\mathbf{V}_{\text{new}})
\end{equation}
where $\mathbf{V}_{\text{new}} \in \mathbb{R}^{N_{\text{new}} \times 384}$. Metadata records are updated with consecutive identifiers starting from $N_{\text{old}}$.

To guarantee crash resilience and prevent concurrent read corruption, index persistence follows a two-stage atomic replacement protocol:
\begin{enumerate}
    \item Vectors and metadata are serialized to temporary disk paths: \texttt{evorag.faiss.tmp} and \texttt{evorag\_meta.json.tmp}.
    \item Upon verified write completion, the operating system executes atomic rename operations via \texttt{os.replace()}, replacing the production index files in a single atomic transaction.
\end{enumerate}

\subsubsection{Local and Accelerated LLM Inference Runtime}
For generative reasoning, EvoRAG supports dual execution configurations:
\begin{itemize}
    \item \textbf{Local Ollama Runtime:} Hosts quantized open-weight instruction models locally via Ollama (\texttt{http://localhost:11434}). The default configuration utilizes \texttt{phi3:mini} (3.8B parameters, 4-bit quantization), which provides high execution throughput on workstation CPUs while maintaining an active RAM footprint under 4.0~GiB. Model weights remain resident between sequential calls by configuring \texttt{keep\_alive = "30m"}.
    \item \textbf{Async Gateway Fallback:} In scenarios where local CPU resources are constrained or when broad general-domain queries exceed local news coverage, an asynchronous HTTPS gateway communicates with Gemini Flash endpoints via persistent \texttt{httpx.AsyncClient} connection pools with token caps and exponential back-off retries.
\end{itemize}

Table~\ref{tab:model_specs} details the core model configurations deployed within the framework.

\begin{table}[H]
\centering
\caption{Model Parameters and Specifications}
\label{tab:model_specs}
\small
\renewcommand{\arraystretch}{1.3}
\begin{tabularx}{\linewidth}{|l|l|l|X|}
\hline
\textbf{Role} & \textbf{Model Name / Tag} & \textbf{Parameters / Dim} & \textbf{Execution Environment} \\
\hline
Text Embedding & \texttt{all-MiniLM-L6-v2} & 22.7M params / 384-D & Local PyTorch / ONNX CPU runtime \\
\hline
Vector Index & FAISS \texttt{IndexFlatL2} & 384-D exact L2 & In-process C++ library wrapper \\
\hline
Local Reasoning & \texttt{phi3:mini} & 3.8B params (4-bit) & Local Ollama server (\texttt{localhost:11434}) \\
\hline
Alternative Local & \texttt{llama3.1:8b} / \texttt{mistral:7b} & 7B--8B params (Q4) & Local Ollama GPU/CPU runtime \\
\hline
Fallback Reasoning & \texttt{gemini-flash-lite-latest} & Cloud MoE Architecture & Asynchronous HTTP/2 gateway \\
\hline
\end{tabularx}
\end{table}

\subsection{Semantic Retrieval and Dynamic Synchronization}
\label{subsec:retrieval_sync}

\subsubsection{Candidate Retrieval and Over-Fetch Deduplication}
When a user submits a query string $Q$, the retrieval engine in \texttt{retriever.py} encodes $Q$ into query vector $\mathbf{q}$. Rather than retrieving only the target $k = 5$ items directly, the system over-fetches a candidate pool of size:
\begin{equation}
k_{\text{cand}} = \min(N, \max(k \cdot M_{\text{mult}}, M_{\text{min}}))
\end{equation}
where candidate multiplier $M_{\text{mult}} = 6$ and minimum pool size $M_{\text{min}} = 20$. 

The candidate chunks are evaluated using document-level identity keys:
\begin{equation}
\mathcal{K}(C) = \{ \text{title}(C), \text{url\_hash}(C), \text{url}(C) \}
\end{equation}
For articles producing multiple top-ranked chunks, only the single chunk with the lowest Euclidean distance is retained. This prevents multiple segments of a single long article from monopolizing the context window, maximizing source diversity across the retrieved context passages.

\subsubsection{Context Sufficiency and Freshness Evaluation}
Before invoking persona reasoning, the retrieved candidates are evaluated by \texttt{evaluate\_retrieval()} across three criteria:
\begin{enumerate}
    \item \textbf{Sufficiency Threshold:} Total word count across retrieved passages must satisfy $\sum \text{words}(C_i) \ge 40$, and passage count must satisfy $|C| \ge 1$.
    \item \textbf{Relevance Threshold:} The best (minimum) L2 Euclidean distance must satisfy:
    \begin{equation}
    d_{\text{min}} = \min_{i} d(\mathbf{q}, \mathbf{v}_i) \le \tau_{\text{rel}}
    \end{equation}
    where $\tau_{\text{rel}} = 1.35$. Queries yielding distances exceeding $1.35$ indicate that the local news corpus contains no semantically related documents.
    \item \textbf{Temporal Freshness:} For queries matching temporal indicators (e.g., \emph{latest, breaking, today, 2026, recent}), publication timestamps are parsed against current UTC time:
    \begin{equation}
    \Delta t_{\text{age}} = t_{\text{current}} - t_{\text{published}} \le 14\text{ days}
    \end{equation}
    If all retrieved documents are older than 14 days for a breaking news query, the retrieval engine flags the context as temporally stale, rerouting the query to broad knowledge synthesis.
\end{enumerate}

\subsubsection{Zero-Restart Dynamic Index Auto-Reloading}
In continuous deployment, the news ingestion scheduler updates the FAISS index on disk every three days. In traditional RAG backends, loading new vectors requires restarting the web application process, causing service downtime. 

EvoRAG implements a zero-restart dynamic index synchronization mechanism in \texttt{retriever.py}. The module tracks the modification timestamp of the index files:
\begin{equation}
T_{\text{disk}} = \max\left(\text{mtime}(\texttt{evorag.faiss}),\ \text{mtime}(\texttt{evorag\_meta.json})\right)
\end{equation}
Upon each incoming request, \texttt{reload\_index\_if\_updated()} compares $T_{\text{disk}}$ against the in-memory timestamp $T_{\text{mem}}$. If $T_{\text{disk}} > T_{\text{mem}}$, the system verifies that the FAISS vector count matches the metadata array length:
\begin{equation}
\text{index.ntotal} \equiv \text{len}(\text{metadata})
\end{equation}
Upon positive validation, the in-memory pointer is updated to point to the newly deserialized index. If a write operation is mid-flight or corrupted, the check fails safely, and the server continues serving queries using the existing valid index without interruption.

\subsection{Multi-Persona Reasoning Engine}
\label{subsec:persona_engine}

\subsubsection{Cognitive Diversity Architecture}
Instead of relying on a single output, EvoRAG employs a parallel generation strategy that produces eight independent analytical responses within a single LLM call. Standard RAG frameworks rely on a single system prompt, which forces the model into a single analytical posture and often mirrors training biases. EvoRAG implements structured cognitive diversity by decomposing analytical reasoning across eight distinct personas. Table~\ref{tab:personas} defines the eight operational personas and their respective reasoning objectives.

\begin{table}[H]
\centering
\caption{The Eight EvoRAG Analytical Personas}
\label{tab:personas}
\small
\renewcommand{\arraystretch}{1.3}
\begin{tabularx}{\linewidth}{|l|p{3.8cm}|X|}
\hline
\textbf{Persona Identifier} & \textbf{Role Title} & \textbf{Core Analytical Function and Directive} \\
\hline
\texttt{analytical-critical} & Analytical-Critical & Deconstructs logical arguments, uncovers hidden assumptions, and evaluates premise validity. \\
\hline
\texttt{empirical-evidential} & Empirical-Evidential & Demands quantitative metrics, verified dates, empirical data points, and strict factual citations. \\
\hline
\texttt{adversarial-skeptical} & Adversarial-Skeptical & Serves as devil's advocate; surfaces failure modes, counter-claims, bias, and conflicting narratives. \\
\hline
\texttt{synthesis-integrative} & Synthesis-Integrative & Bridges disparate facts, harmonizes conflicting passages, and identifies overarching thematic patterns. \\
\hline
\texttt{domain-expert} & Domain-Expert & Explains complex domain mechanics with precise terminology, contextual depth, and theoretical rigor. \\
\hline
\texttt{temporal-contextual} & Temporal-Contextual & Evaluates historical antecedents, timelines, rate of development, and forward trajectories. \\
\hline
\texttt{risk-evaluative} & Risk-Evaluative & Identifies systemic vulnerabilities, threat vectors, regulatory exposure, and impact severity. \\
\hline
\texttt{creative-lateral} & Creative-Lateral & Draws lateral analogies, examines interdisciplinary impacts, and explores non-obvious scenarios. \\
\hline
\end{tabularx}
\end{table}

\subsubsection{Prompt Engineering and Context Citation Assembly}
Each persona prompt combines the persona's specialized system directive with the shared contextual context passages assembled from the retrieved chunks. Context blocks are cited systematically:
\begin{equation}
\label{eq:context_block}
\text{Block}_i = \text{``[}i\text{] Source: } S_i \text{ | Title: } T_i \text{ $\backslash$n Text: } P_i\text{''}
\end{equation}

\subsubsection{Asyncio Semaphore Concurrency Control}
To execute eight independent persona analyses on resource-constrained development hardware without exhausting CPU or VRAM bandwidth, \texttt{personas.py} manages invocation concurrency using an \texttt{asyncio.Semaphore}:
\begin{equation}
\Sigma = \text{asyncio.Semaphore}(C_{\text{max}})
\end{equation}
where $C_{\text{max}} = 1$ by default for local CPU execution, ensuring serialized Ollama execution that prevents thread contention, memory thrashing, and socket timeouts. All calls are scheduled concurrently using \texttt{asyncio.gather()}, allowing personas to execute as threads become available.

\subsection{Background Knowledge Ingestion and Orchestration}
\label{subsec:scheduler_orchestration}
In \texttt{knowledge\_scheduler.py}, background knowledge base maintenance is handled using Python's \texttt{APScheduler} library via a \texttt{BlockingScheduler} configured on a 3-day interval:
\begin{equation}
\text{Trigger: } \texttt{IntervalTrigger(days=3)}
\end{equation}

When the scheduler executes, it queries the status file \texttt{index/ingestion\_status.json}. If the elapsed time since the last successful ingestion satisfies:
\begin{equation}
\Delta t_{\text{update}} = t_{\text{current}} - t_{\text{last\_success}} < 3\text{ days}
\end{equation}
the scheduler logs an index-reuse confirmation and sleeps until the next scheduled interval, avoiding redundant network API requests and CPU utilization.

\subsection{Web Client, Multi-Turn Chat, and Voice Modalities}
\label{subsec:web_interface}

\subsubsection{FastAPI REST Architecture and Service Lifespan}
The web application service is implemented in \texttt{api.py} using FastAPI and Uvicorn. The service exposes CORS-enabled JSON endpoints alongside static asset serving. The lifecycle context manager (\texttt{lifespan}) initializes the SQLite database tables on startup, verifies the FAISS index, and handles clean resource deallocation on shutdown.

\subsubsection{Multi-Turn Conversation Persistence and User Isolation}
To support longitudinal multi-session research, \texttt{conversations.py} implements conversation state management. Each authenticated user owns isolated conversation threads tracked in SQLite:
\begin{itemize}
    \item \textbf{Chat Session Auto-Naming:} The \texttt{clean\_title()} utility parses the first user query, stripping conversational lead-in prefixes (e.g., \emph{``what is'', ``explain to me'', ``can you tell me about''}), sanitizing punctuation, and truncating to a readable 40-character title.
    \item \textbf{Chronological Thread Retrieval:} Previous conversations are loaded through \texttt{GET /conversations/\{id\}/messages}, restoring all user questions and consensus answers in sequence without crosstalk across user accounts.
\end{itemize}

\subsubsection{Browser-Native Speech-to-Text (STT)}
Voice input is implemented directly in \texttt{frontend/index.html} using the browser's native Web Speech API (\texttt{webkitSpeechRecognition} / \texttt{SpeechRecognition}). The implementation avoids third-party paid audio APIs, operating locally within the browser engine:
\begin{itemize}
    \item \textbf{Continuous Dictation:} The recognizer operates with \texttt{continuous = false} and \texttt{interimResults = true}, displaying transcribed tokens inside the query textarea in real time.
    \item \textbf{Multilingual Recognition:} A UI language selector allows users to toggle recognition across Indian regional and international languages: English (\texttt{en-IN}), Kannada (\texttt{kn-IN}), Hindi (\texttt{hi-IN}), Tamil (\texttt{ta-IN}), and Telugu (\texttt{te-IN}).
    \item \textbf{Non-Destructive Review:} Speech input is populated directly into the textarea, allowing the user to inspect and edit recognized text before clicking Submit.
\end{itemize}

\subsubsection{Browser-Native Text-to-Speech (TTS)}
Spoken response synthesis utilizes the browser's native \texttt{speechSynthesis} interface. Each consensus answer card features audio playback controls (Listen, Pause, Resume, Stop). The synthesis engine filters the operating system's installed voice inventory, selects a native voice matching the active language preference, and manages playback states across rapid user clicks to avoid overlapping audio streams.

\clearpage

\section{Core Algorithms}

The technical implementation of EvoRAG relies on several core algorithmic concepts,
with the primary orchestration handled by the EvoRAG Pipeline described below.

\subsubsection{Evolutionary Multi-Persona RAG Algorithm}

\vspace{0.5em}

\begin{algorithm}[H]
\caption{EvoRAG Pipeline --- Query to Consensus Response}
\label{alg:query_consensus}
\small
\SetAlgoLined
\KwIn{$user\_query$}
\KwOut{$best\_persona,\ best\_answer,\ score\_breakdown,\ sources$}

\BlankLine
\tcp{Step 1: Retrieval}
$chunks \leftarrow \text{FAISS\_Retrieve}(user\_query,\ top\_k=3)$\;
\If{$chunks = \emptyset$}{
    \Return \textit{``No relevant context found.''}\;
}
$context \leftarrow \text{Merge}(chunks)$\;

\BlankLine
\tcp{Step 2: Multi-Persona Generation}
$personas \leftarrow \{$Analytical-Critical, Empirical-Evidential, Adversarial-Sceptical,
Creative-Lateral, Domain-Expert, Temporal-Contextual, Risk-Evaluative, Synthesis-Integrative$\}$\;
$prompt_{gen} \leftarrow \text{BuildPrompt}(user\_query,\ context,\ personas)$\;
$drafts \leftarrow \text{ParseJSON}(\text{LLM}(prompt_{gen}))$\;

\BlankLine
\tcp{Step 3: Peer-Review Scoring}
$prompt_{judge} \leftarrow \text{BuildJudgePrompt}(user\_query,\ drafts)$\;
$scores \leftarrow \text{ParseJSON}(\text{LLM}(prompt_{judge}))$\;
\tcp*[r]{Factual Grounding (0--10) + Analytical Depth (0--10) + Clarity (0--10) $= /30$}

\BlankLine
\tcp{Step 4: Consensus Selection}
$best\_persona \leftarrow \arg\max_{p \in personas}\ scores[p]$\;
$best\_answer \leftarrow drafts[best\_persona].answer$\;

\BlankLine
\tcp{Step 5: Output}
\Return $\{best\_persona,\ best\_answer,\ scores,\ chunks.metadata\}$\;
\end{algorithm}

\vspace{1em}

\begin{algorithm}[H]
\caption{Automated News Ingestion and Incremental Index Update}
\label{alg:scheduler_pipeline}
\small
\SetAlgoLined
\KwIn{Query topics $\mathcal{Q}$, update interval $T_{\text{interval}} = 3\text{ days}$}
\KwOut{Updated FAISS index $\mathcal{I}$, updated metadata $\mathcal{M}$, status log}

\BlankLine
$status \leftarrow \text{LoadStatus}(\texttt{index/ingestion\_status.json})$\;
\If{$\text{IndexExists}()$ \textbf{and} $(t_{\text{now}} - status.\text{last\_update} < T_{\text{interval}})$}{
    \Return \textit{``Index fresh. Reusing existing FAISS vectors.''}\;
}

\BlankLine
\tcp{Step 1: Multi-Source Harvesting}
$chunks_{\text{raw}} \leftarrow \emptyset$\;
\ForEach{$q \in \mathcal{Q}$}{
    $chunks_{\text{raw}} \leftarrow chunks_{\text{raw}} \cup \text{FetchNewsAPI}(q) \cup \text{FetchGNews}(q) \cup \text{FetchRSS}(q)$\;
}

\BlankLine
\tcp{Step 2: Cleaning and Signature Filtering}
$\mathcal{S}_{\text{existing}} \leftarrow \text{ExtractSignatures}(\mathcal{M})$\;
$chunks_{\text{new}} \leftarrow \emptyset$\;
\ForEach{$c \in chunks_{\text{raw}}$}{
    $c_{\text{clean}} \leftarrow \text{CleanAndSanitize}(c)$\;
    $sig \leftarrow (\text{MD5}(c_{\text{clean}}.\text{url}),\ \text{SHA256}(c_{\text{clean}}.\text{text}))$\;
    \If{$sig \notin \mathcal{S}_{\text{existing}}$}{
        $chunks_{\text{new}} \leftarrow chunks_{\text{new}} \cup \{c_{\text{clean}}\}$\;
        $\mathcal{S}_{\text{existing}} \leftarrow \mathcal{S}_{\text{existing}} \cup \{sig\}$\;
    }
}

\BlankLine
\tcp{Step 3: Incremental Embedding and Atomic Commit}
\If{$chunks_{\text{new}} \neq \emptyset$}{
    $\mathbf{V}_{\text{new}} \leftarrow \text{MiniLM\_Encode}(chunks_{\text{new}}.\text{text},\ \text{batch\_size}=64)$\;
    $\mathcal{I}.\text{add}(\mathbf{V}_{\text{new}})$\;
    $\mathcal{M}.\text{append}(chunks_{\text{new}}.\text{metadata})$\;

    $\text{WriteAtomic}(\texttt{evorag.faiss.tmp}, \mathcal{I})$\;
    $\text{WriteAtomic}(\texttt{evorag\_meta.json.tmp}, \mathcal{M})$\;
    $\text{AtomicSwap}(\texttt{.tmp} \rightarrow \texttt{production})$\;
}

$status.\text{last\_update} \leftarrow t_{\text{now}}$\;
$status.\text{vector\_count} \leftarrow \mathcal{I}.\text{ntotal}$\;
$\text{SaveStatus}(\texttt{index/ingestion\_status.json}, status)$\;
\Return $status$\;
\end{algorithm}

\section{Consensus-Based Evaluation}
\label{sec:consensus_evaluation}

The system implements a self-regulating peer-review mechanism.

\begin{itemize}
    \item \textbf{Peer Voting Mechanism:} Each persona evaluates all generated responses, acting as a “judge” for others. Self-scoring is strictly excluded ($i = j$) to prevent bias.
    
    \item \textbf{Scoring Criteria:} Responses are graded on three dimensions: Factual Grounding ($s^{\text{fact}} \in [1, 10]$), Analytical Depth ($s^{\text{reas}} \in [1, 10]$), and Clarity ($s^{\text{comp}} \in [1, 10]$). The composite pairwise score awarded by persona $p_i$ to peer $p_j$ is:
    \begin{equation}
    S_{i,j} = \frac{1}{3} \left( s_{i,j}^{\text{fact}} + s_{i,j}^{\text{reas}} + s_{i,j}^{\text{comp}} \right), \quad \forall i \neq j
    \end{equation}
    
    \item \textbf{Final Selection:} The unweighted consensus score for persona $p_j$ is obtained by averaging across all peer evaluations:
    \begin{equation}
    \bar{S}_j = \frac{1}{N-1} \sum_{i \neq j} S_{i,j}
    \end{equation}
    The highest-scoring response, representing the consensus of the multi-persona reasoning session, is chosen as the final output alongside its source citations and peer matrix.
\end{itemize}

\section{Evolutionary Learning Engine}
\label{sec:evolutionary_engine}

To ensure continuous improvement without retraining the base LLM, an evolutionary layer is added.

\begin{itemize}
    \item \textbf{Performance Tracking:} The win rates and qualitative scores of each persona are stored in a local SQLite database and serialized to \texttt{score\_store.json}. Rolling averages are maintained over a window of $N_w = 30$ queries:
    \begin{equation}
    \mu_p = \frac{1}{|\mathcal{H}_p|} \sum_{q \in \mathcal{H}_p} \bar{S}_{p,q}, \quad \rho_p = \frac{\sum_{q \in \mathcal{H}_p} \mathbb{I}(p \in \mathcal{P}_{\text{top}, q})}{|\mathcal{H}_p|}
    \end{equation}
    
    \item \textbf{Adaptive Replacement:} Personas that consistently fail to win the peer-review or receive poor user feedback are periodically replaced with new persona variants. A persona is flagged for replacement if it has participated in at least $N_{\text{min}} = 10$ queries and its rolling mean score drops below threshold $\tau_{\text{replace}} = 7.0$:
    \begin{equation}
    \mu_p < \tau_{\text{replace}}
    \end{equation}
    
    \item \textbf{Self-Improvement:} This allows the system to evolve its intellectual makeup over time based on the specific types of queries it receives, recording prompt mutation lineages in \texttt{versions.json} without requiring service restarts.
\end{itemize}

\section{User Adaptation and Personalization}
\label{sec:user_personalization}

The system adapts to individual user needs through a feedback loop.

\begin{itemize}
    \item \textbf{Feedback Integration:} The UI records user likes and dislikes for each response, as well as manual persona selection overrides from the interactive carousel.
    
    \item \textbf{Affinity Modeling:} Personalized preference profiles are built, identifying which personas or response styles the user prefers. The net user affinity ratio is calculated as:
    \begin{equation}
    \alpha(u, p) = \frac{N_{\text{pos}}(u, p) - N_{\text{neg}}(u, p)}{N_{\text{total}}(u, p) + \epsilon}
    \end{equation}
    where $\epsilon = 1.0$ serves as a Laplacian smoothing factor for sparse user interactions.
    
    \item \textbf{Adaptive Output:} The voting weights in the consensus phase are dynamically adjusted based on these preference profiles:
    \begin{equation}
    \tilde{S}_j(u) = \bar{S}_j \cdot \left( 1 + \gamma \cdot \alpha(u, p_j) \right)
    \end{equation}
    where $\gamma = 0.20$ bounds user preference influence to $\pm 20\%$, harmonizing subjective user preference with objective peer-review consensus.
\end{itemize}
